\chapter{路径积分}

\section{量子力学回顾}

近代光学实验中，光常同时表现出粒子和波动的性质，人们逐渐形成波粒二象性的观念，笼统地称为\textbf{量子态}(quantum state)。
量子一词起初源于黑体辐射的能量离散化，后来延伸为波粒二象性等有悖于经典力学的性质。
其中一个有趣的事实是，我们无法同时观测粒子的具体位置及动量 $\bm x,\bm p$。
Heisenberg 提出，经典的测量值 $\bm x,\bm p$ 是实数，但二者的不确定性意味着必须将其“提升”为一种\textbf{算符}(operator) $\hat{\bm x},\hat{\bm p}$，作用于物质的量子态上。
考虑一维情形，他根据实验和已有研究提炼出
\eq{
    [\hat x,\hat p]:=\hat x\hat p - \hat p\hat x = \i\hbar,
}
$\hbar=h/2\pi$ 是约化 Planck 常数。这称为 $\hat x,\hat p$ 的\textbf{对易子}，意思无非是对量子态先作用 $\hat p$ 再作用 $\hat x$，减去其相反操作，结果相当于乘上 $\i\hbar$。所以量子力学不可避免涉及复数。
尽管我们仍没有说出量子态、算符的数学表示，但我们要求它们一定服从这种对易关系。
经典力学是 $h\to 0$ 的极限，意味着 $\hat x\hat p=\hat p\hat x$，作为普通实数的 $x,p$ 就满足这个可交换乘法。
算符的非对易性可以说是量子性质的关键，因此这种“算符提升”也称为从经典力学走向量子力学的\textbf{量子化}。

下面我们来猜测算符可能的数学表示。人们现已广泛清楚，矩阵的乘法是不对易的，因此矩阵是合适的数学候选。如此，量子态就是列矢量，通常写作 $\ket{}$。一般会在里面放一些符号，仅表标签含义，比如 $\ket{\psi}$。
但在当时，矩阵并非必学知识。Heisenberg 从经典力学的 Poisson 括号绕了一圈才回到了矩阵，相当于重新发明了矩阵。可以说，今日之大学得以普及线性代数，有 Heisenberg 的功劳。
不过这里会比我们学的线性代数复杂：因为可以证明，任意有限维度的矩阵都无法凑出对易关系。因此，算符必须是无限维空间上的矩阵。
相较于 Heisenberg，Schr\"odinger 更关注于量子态的具体表示，方法会更简单些。
他考虑的算符是直接的偏微分算符，比如 $\pdv{x},\pdv{t}$。因此量子态是某种实变但复值的函数 $\psi(x,t)$，称为\textbf{波函数}。
两种观点在后来被证明为等价的。
这也解释了为什么所涉及的线性代数一定是无穷维的，因为函数空间是无穷维的。



关键在于物理量的测量。由于现在提升为算符，可能的测量值只能解释为算符的某个本征值。如果 $\hat A$ 具有分立的本征值 $\{a_i\}$，则测量 $A$ 的结果只能是 $a_i$ 中的一个。回忆一下，本征值的定义无非是 $\hat A u_n=a_n u_n$，对应的 $u_n$ 称为 $\hat A$ 的本征函数。
Born 根据杨⽒双缝实验提出，测量结果 $a_i$ 出现的概率由态决定：$\psi=\sum_n c_n u_n$，则相应概率为 $|c_n|^2$。连续版本为



比如位置测量值 $x$ 是位置算符 $\hat x$ 的本征值：$\hat x\ket{x}=x\ket{x}$。$\ket{x}$ 是 $\hat x$ 的本征矢，所以用了标签 $x$，仅此而已。当然，两种观点分别对应的本征矢、本征函数可统称为\textbf{本征态}。

$|\langle a_i|\psi\rangle|^2$，其中 $\ket{\psi}$ 是系统的量子态。






：将 $|\psi|^2$ 解释为 $t$ 时刻粒子出现在 $\bm x$ 的概率密度分布。
于是通常默认归一化条件：
\eq{
    \int\d[3]{x} |\psi|^2 = 1.
}
位置平均值为 $\langle x \rangle=\int\d[3]{x} x|\psi|^2 =\int\d[3]{x} \psi^* x\psi$。算符的平均值为 $\langle A \rangle = \int\d[3]{x} \psi^* \hat A \psi$。位置算符的平均值为位置平均值，则必须 $\hat x = x$。
由 Leibniz 乘法法则得 $\hat p : = -\i\hbar\pdv{x}$ 满足对易关系。






通过波动力学的一些类比方法，Schr\"odinger 指出波函数应满足方程
\eq{
    \i\pdv{t}\psi=\left(-\frac{\hbar^2}{2m}\Delta+V(\bm x,t)\right)\psi.
}
可见若 $\{\psi_i\}$ 是解，则任意\textbf{线性叠加} $\sum_i c_i\psi_i$ 也是解。
这也意味着 $\psi$ 的重标度 $c\psi$ 仍是解，故总可以对波函数归一化。
$V$ 为常数时，上式具有简单的平面波解 $\psi=A\e^{\i(\bm k\cdot\bm x - \omega t)}$。容易发现 $\bm k,\omega$ 必须满足 $\omega = \hbar \bm k^2/2m + V/\hbar$。借 Planck 黑体辐射、Einstein 光电效应，de Broglie 指出，粒子的能量、动量可与波函数的波动性质相联系：
\eq{
    E = h\nu = \hbar\omega,\quad \bm p = \hbar\bm k.
}
因此这无非是说粒子满足经典力学而非相对论的关系式
\eq{
    E = {\bm p^2}/{2m} + V.
}


容易发现任意 $V$ 仍然具有平面波解。
简单平面波的线性叠加也是解。

自由单粒子的叠加


$[-L/2,L/2]$ 上良好的函数可展开为 Fourier 级数
\[f(x)=\sum_{N=-\infty}^\infty \frac{F(p_N)}{L} \e^{\i p_N x},\quad F(p_N)=\int_{-L/2}^{L/2}\d{x}\e^{-\i p_N x}f(x),\]
其中 $p_N={2\pi N}/{L}$ 称为\textbf{模式}(mode)，$F(p)$ 是 $f(x)$ 的 Fourier 变换。令 $L\to\infty$，则 $\Delta p=2\pi/L\to 0$，级数化为积分，即 $f(x)=\int \frac{\d p}{2\pi} \e^{\i p x}F(p),F(p)=\int\d{x}\e^{-\i p x}f(x)$。任意 $f(x)$ 以 $\e^{\i px}$ 为基底。
同理，取 $V=L^n$ 并连续化，可得
\[f(x)=\int \frac{\d[n]p}{(2\pi)^n} \e^{\i p_\mu x^\mu}F(p),\quad F(p)=\int\d[n]{x}\e^{-\i p_\mu x^\mu}f(x).\]


当有多个粒子时，比如两个粒子，系统的量子态 $\psi(\bm x_1,\bm x_2,t)$ 是两个位置变量的函数。对于 $N$ 个粒子，系统的量子态 $\psi(\bm x_1,\cdots,\bm x_N,t)$ 是 $3N$ 个位置变量的函数。因此，粒子数的增加导致系统的量子态空间维数呈指数级增长。

波函数不再是单个粒子的位置的函数，而是所有粒子位置的函数。这种描述方式虽然在理论上是可行的，但在实际计算中非常复杂，尤其是当粒子数目较多时。
引入态矢量的概念后，系统的量子态可以表示为 $\ket{\psi}$，而不是具体的波函数形式。这种抽象的表示方式使得我们能够更灵活地处理多粒子系统，尤其是在涉及到粒子交换对称性和量子统计时。


因此，描述多粒子系统的量子态空间非常庞大，难以直接处理。为了解决这个问题，物理学家引入了\textbf{量子场}的概念，将粒子视为量子场的激发态，从而能够更有效地描述多粒子系统的行为。这种方法不仅简化了多粒子系统的描述，还为理解粒子之间的相互作用提供了一个更自然的框架。



\section{路径积分方法}

我们曾在经典力学中学习过最小作用量原理。Feynman 对此有一个进一步的解释。
